\chapter*{Introduction Générale}
\addcontentsline{toc}{chapter}{Introduction Générale}

L'avènement des technologies de l'information et de la communication a profondément bouleversé le paysage éducatif mondial, induisant une transition fulgurante des modèles pédagogiques traditionnels vers des environnements d'apprentissage numériques et hybrides. Cette évolution a mis en exergue l'importance capitale des systèmes de gestion de l'apprentissage (LMS), devenus les pierres angulaires de la numérisation des processus éducatifs.

Cependant, les LMS conventionnels souffrent de limitations intrinsèques. L'un des défis majeurs réside dans leur approche pédagogique statique, proposant un parcours identique à tous les étudiants, sans considération pour leurs acquis préalables ou leur rythme d'assimilation. Parallèlement, l'essor de l'enseignement à distance a soulevé des préoccupations majeures quant à l'intégrité académique et à la fiabilité des évaluations. L'éloignement physique ouvre la porte à diverses formes de fraudes, rendant impératif le déploiement de mécanismes de télésurveillance (proctoring).

C'est à la confluence de ces deux problématiques que se situe le projet \textbf{Fraudly}. L'objectif principal de ce travail est de concevoir et développer une plateforme d'apprentissage intelligente de nouvelle génération : un \textit{Agentic Learning Management System}. L'originalité réside dans l'intégration de l'intelligence artificielle générative et analytique au cœur du processus éducatif. Des agents autonomes ont pour mission de personnaliser les parcours, d'assister les étudiants via un tutorat contextuel, de générer des évaluations complexes, et d'assurer une télésurveillance multimodale et comportementale.

Ce présent manuscrit est structuré en sept chapitres détaillant l'ensemble de notre démarche d'ingénierie logicielle. Le premier chapitre explore le contexte général et la problématique. Le deuxième chapitre analyse en profondeur les besoins et cas d'utilisation de la plateforme. Le troisième chapitre se consacre à l'architecture globale en microservices. Les chapitres quatre, cinq et six détaillent la conception fonctionnelle, l'intégration de l'intelligence artificielle et le module de proctoring. Enfin, le septième chapitre présente l'implémentation technologique avant de clore sur une conclusion générale.


Nous espérons que ce document mettra en évidence l'étendue de nos recherches, la complexité de l'architecture mise en œuvre et l'impact potentiel de l'intelligence artificielle agentique sur l'avenir de la pédagogie numérisée. En vous souhaitant une bonne lecture.